\begin{Lemma}\label{lem:Schur}Assume that $\F$ is algebraically closed. If $V$ is a finite-dimensional irreducible $\Re$-module, then each central element of~$\Re$ acts on~$V$ as scalar multiplication.
\end{Lemma}
\begin{proof}This result follows from applying Schur's lemma to $\Re$.
\end{proof}

\begin{Lemma}\label{lem:theta}For any $i\in \Z$, each of the following hold:
\begin{enumerate}\itemsep=0pt
\item[$(i)$] $\theta_{i+1}+\theta_{i-1}=2(\theta_i+1)$,
\item[$(ii)$] $\theta_{i+1}\theta_{i-1}=\theta_i(\theta_i-2)$.
\end{enumerate}
\end{Lemma}
\begin{proof}The result can be routinely verified using (\ref{theta_i}).
\end{proof}

\begin{Theorem}\label{thm:surjective}Assume that $\F$ is algebraically closed with $\operatorname{char}\F=0$. Let $d$ denote a nonnegative integer. If $V$ is a $(d + 1)$-dimensional irreducible $\Re$-module, then there exist $a, b, c \in \F$ such that the $\Re$-module $R_d(a, b, c)$ is isomorphic to~$V$.
\end{Theorem}
\begin{proof}Given any scalar $\kappa \in \F$, we define
\begin{gather*}
\vartheta_i(\kappa)
=(\kappa-i)(\kappa-i+1)
\qquad
\hbox{for all $i\in \Z$}.
\end{gather*}
Since $\operatorname{char}\F=0$, for any distinct integers $i,j$, the scalars $\vartheta_i(\kappa)$ and $\vartheta_j(\kappa)$ are equal if and only if $i+j=2\kappa+1$. In particular $\{\vartheta_i(\kappa)\}_{i=0}^{-\infty}$ contains infinitely many values.

Since $\F$ is algebraically closed, we may choose a scalar $\kappa\in \F$ such that $\vartheta_0(\kappa)$ is an eigenvalue of~$A$ on~$V$. Since~$V$ is of dimension $d+1$, there are at most $d+1$ distinct eigenvalues of~$A$ on~$V$. Thus, there exists an integer $j\leq 0$ such that $\vartheta_j(\kappa)$ is an eigenvalue of~$A$ but $\vartheta_{j-1}(\kappa)$ is not an eigenvalue of~$A$ on~$V$. Set
\begin{gather*}
a=\kappa-j-\tfrac{d}{2}.
\end{gather*}
Similarly, there exists a scalar $\lambda\in \F$ and an integer $k\leq 0$ such that $\vartheta_k(\lambda)$ is an eigenvalue of~$B$ but $\vartheta_{k-1}(\lambda)$ is not an eigenvalue of~$B$ in~$V$. We set
\begin{gather*}
b=\lambda-k-\tfrac{d}{2}.
\end{gather*}
Observe that under these settings, we have
\begin{gather}
\theta_i = \vartheta_{i+j}(\kappa) \qquad\hbox{for all $i\in \Z$},\label{setting:a}\\
\theta_i^* = \vartheta_{i+k}(\lambda)\qquad\hbox{for all $i\in \Z$}.\label{setting:b}
\end{gather}
By Lemma \ref{lem:Schur}, the element $\delta$ acts on $V$ as scalar multiplication. Since $\F$ is algebraically closed, there exists a scalar $c\in \F$ such that the action of $\delta$ on $V$ is the scalar multiplication by
\begin{gather*}
\eta=\textstyle\frac{d}{2}\big(\frac{d}{2}+1\big)+a(a+1)+b(b+1)+c(c+1).
\end{gather*}
To prove the theorem, it now suffices to show that there exists an $\Re$-module isomorphism from $R_d(a,b,c)$ into~$V$.

Given any $T\in \Re$ and $\theta\in \F$, we let
\begin{gather*}
V_T(\theta)=\{v\in V\,|\, Tv=\theta v\}.
\end{gather*}
Pick any $v\in V_A(\theta_0)$. Applying each side of~(\ref{AAB}) to~$v$ and using Lemma~\ref{lem:theta} to simplify the result, we obtain that
\begin{gather}\label{AAB_V(theta)}
(A-\theta_{-1})(A-\theta_1) B v=2(\theta_0(\theta_0-\eta)+\alpha)v.
\end{gather}
Left multiplying each side of (\ref{AAB_V(theta)}) by $(A-\theta_0)$, we obtain that
\begin{gather*}
(A-\theta_{-1})(A-\theta_0)(A-\theta_1) B v=0.
\end{gather*}
By (\ref{setting:a}), the scalar $\theta_{-1}$ is not an eigenvalue of $A$ in $V$. Hence
\begin{gather*}
(A-\theta_0)(A-\theta_1) B v=0.
\end{gather*}
In other words $(A-\theta_1) B v\in V_A(\theta_0)$ and therefore $V_A(\theta_0)$ is invariant under $(A-\theta_1) B$. Since~$\F$ is algebraically closed, there exists an eigenvector $u$ of $(A-\theta_1) B$ in $V_A(\theta_0)$. Similarly, there exists an eigenvector $w$ of $(B-\theta_1^*)A$ in $V_B(\theta_0^*)$. Define
\begin{gather}
u_i=\prod_{h=0}^{i-1} (B-\theta_h^*) u \qquad \hbox{for all $i\in \N$},\label{vi}\\
w_i=\prod_{h=0}^{i-1} (A-\theta_h) w\qquad \hbox{for all $i\in \N$}. \label{wi}
\end{gather}

We now proceed by induction to show that
\begin{gather}\label{claim}
(A-\theta_i)u_i\in \operatorname{span}_\F\{u_0,u_1,\ldots,u_{i-1}\} \qquad \hbox{for all $i\in \N$}.
\end{gather}
Since $u$ is an eigenvector of $(A-\theta_1) B$ in $V_A(\theta_0)$, the claim is true for $i=0,1$.
Now suppose that $i\geq 2$. Applying each side of~(\ref{ABB}) to $u_{i-2}$, we obtain that
\begin{gather}\label{ABBvi-2}
\big(A B^2 -2 BAB +B^2 A - 2 A B-2 BA-2 B^2+2 \eta B\big) u_{i-2}=-2 \beta u_{i-2}.
\end{gather}
By Lemma \ref{lem:Schur}, the right-hand side of (\ref{ABBvi-2}) is a scalar multiple of~$u_{i-2}$. Using the inductive hypothesis, (\ref{vi}), and Lemma~\ref{lem:theta}(i), we find that the left-hand side of~(\ref{ABBvi-2}) is equal to
\begin{gather}\label{ABBvi-2:LH'}
(A-\theta_i) u_i
\end{gather}
plus an $\F$-linear combination of $u_0,u_1,\ldots,u_{i-1}$. Combining the above results, the claim~(\ref{claim}) follows.

Next, we show that $\{u_i\}_{i=0}^d$ is an $\F$-basis for $V$. Suppose on the contrary that there is an integer $h$ with $0\leq h\leq d-1$ such that $u_{h+1}$ is an $\F$-linear combination of $u_0,u_1,\ldots,u_h$. Let~$W$ denote the $\F$-subspace of~$V$ spanned by $u_0,u_1,\ldots,u_h$.
Since~$W$ is $B$-invariant by~(\ref{vi}) and $A$-invariant by~(\ref{claim}), it follows that $W$ is an $\Re$-submodule of $V$ by Lemma~\ref{lem:delta}(ii). Since~$V$ is irreducible, this forces that $W=V$.
By construction, $W$ is of dimension at most~$d$, a~contradiction. Therefore $\{u_i\}_{i=0}^d$ is an $\F$-basis for $V$.
By a similar argument, it follows that
\begin{gather}\label{claim'}
(B-\theta_i^*) w_i\in \operatorname{span}_\F\{w_0,w_1,\ldots,w_{i-1}\}\qquad \hbox{for all $i\in \N$}
\end{gather}
and $\{w_i\}_{i=0}^d$ is an $\F$-basis for $V$.

By (\ref{claim}), the matrix representing $A$ with respect to the $\F$-basis $\{u_i\}_{i=0}^d$ is upper triangular with diagonal entries $\{\theta_i\}_{i=0}^d$. By the Cayley--Hamilton theorem, we have
\begin{gather}\label{CH}
\prod\limits_{i=0}^d (A-\theta_i) w_0=0.
\end{gather}
In other words $Aw_d=\theta_d w_d$ by (\ref{wi}). Hence the matrix representing $A$ with respect to the $\F$-basis $\{w_i\}_{i=0}^d$ is
\begin{gather*}
\begin{pmatrix}
\theta_0 & & & &{\bf 0}
\\
1 &\theta_1
\\
&1 &\theta_2
 \\
& &\ddots &\ddots
 \\
{\bf 0} & & &1 &\theta_d
\end{pmatrix}.
\end{gather*}

By (\ref{claim'}), the matrix representing $B$ with respect to the $\F$-basis $\{w_i\}_{i=0}^d$ is upper triangular with diagonal entries $\{\theta_i^*\}_{i=0}^d$. We let $\{\varphi_i'\}_{i=1}^d$ denote its superdiagonal entries as follows
\begin{gather}\label{B}
\begin{pmatrix}
\theta_0^* &\varphi_1' & & &{\bf *}
\\
 &\theta_1^* &\varphi_2'
\\
 & &\theta_2^* &\ddots
 \\
 & & &\ddots &\varphi_d'
 \\
{\bf 0} & & & &\theta_d^*
\end{pmatrix}.
\end{gather}
Applying each side of (\ref{AAB}) to $w_{i-1}$, we obtain from the coefficients of $w_i$ that
\begin{gather}\label{varphi'}
\varphi_{i+1}'-2\varphi_i'+\varphi_{i-1}'
=(\theta_i-\theta_{i-1}+2)\theta_i^*
-(\theta_i-\theta_{i-1}-2)\theta_{i-1}^*
+2(\theta_i+\theta_{i-1}-\eta)
\end{gather}
for all $1\leq i\leq d$, where $\varphi_0'$ and $\varphi_{d+1}'$ are interpreted as zero.
It is straightforward to verify that $\{\varphi_i\}_{i=1}^d$ also satisfy the recurrence relation (\ref{varphi'}). Since $\operatorname{char}\F=0$, the corresponding homogeneous recurrence relation
\begin{gather*}
\sigma_{i+1}-2\sigma_i+\sigma_{i-1}=0, \qquad 1\leq i\leq d,
\end{gather*}
with the initial values $\sigma_0=0$ and $\sigma_{d+1}=0$, has the unique solution $\sigma_i=0$ for all $0\leq i\leq d+1$. Therefore $\varphi_i'=\varphi_i$ for all $1\leq i\leq d$.

Up to this point, we have shown that
\begin{gather}
Bw_0=\theta_0^* w_0,\label{even:U1}\\
(B-\theta_1^*)(A-\theta_0)w_0=\varphi_1 w_0,\label{even:U2}\\
\delta w_0=\eta w_0.\label{even:U3}
\end{gather}

Applying each side of (\ref{ABB}) to $w_0$ and using~(\ref{even:U2}) to simplify the resulting equation, we find that
\begin{gather}\label{even:U4}
\beta w_0=\zeta^* w_0.
\end{gather}
Using logic similar to what was used to show (\ref{even:U2}), we obtain $(A-\theta_1)(B-\theta_0^*)u_0=\varphi_1 u_0$. Now, by applying each side of (\ref{AAB}) to $u_0$ and using the above equation to simplify the resulting equation, we find that $\alpha u_0=\zeta u_0$. It follows from Lemma~\ref{lem:Schur} that
\begin{gather}\label{even:U5}
\alpha w_0=\zeta w_0.
\end{gather}
In view of (\ref{even:U1})--(\ref{even:U5}), it follows from Proposition~\ref{prop:universal} that there exists a unique $\Re$-module homomorphism $M_d(a,b,c)\to V$ that sends~$m_0$ to~$w_0$. By Proposition~\ref{prop:Verma}(i), the entries above the superdiagonal in~(\ref{B}) are zero.
Combining the above $\Re$-module homomorphism \mbox{$M_d(a,b,c)\to V$} with~(\ref{CH}), there is an $\Re$-module homomorphism
\begin{gather}\label{R->V}
R_d(a,b,c)\to V
\end{gather}
that sends $v_0$ to $w_0$ by Proposition~\ref{prop:R}. Since the $\Re$-module $V$ is irreducible, the homomorphism~(\ref{R->V}) is onto. Since $R_d(a,b,c)$ and $V$ are both of dimension $d+1$, it follows that (\ref{R->V}) is an isomorphism. The result follows.
\end{proof}
